
\chapter{Dates and times}

\index{time@``time''}

One of the most common categories of information you'll ever deal with are
those dealing with various aspects of \textit{time}. It turns out this is
considerably more complex than you might think at first. For one thing, the
English word ``time'' can mean several different things -- a point in time
(12:05pm on Feb.~20, 2020), an abstract wall clock time (12:05pm), a duration
of time (1 hour and 15 minutes), and possibly other things.

For another, the clock and calendar systems that humans have devised are just
plain weird. Why does the break between ``am'' and ``pm'' occur between 11 and
12 (11am, 12\underline{\textbf{p}}m, 1pm) instead of between 12 and 1 where any
sane person would put it? Why do some months have more days than others? Heck,
they don't even follow the moon. And why do some \textit{years} have more days
than others? And why doesn't every month start over on a Monday? And so forth.
There are good answers to some of these questions, and lousy answers to others,
but they all make dealing with time a major pain.

\section{The \texttt{datetime} package}

\index{import@\texttt{import}}
\index{datetime@\texttt{datetime} package}

Python helps deal with the pain through the \texttt{datetime} package. You have
to import it, of course, which I usually do via:

\begin{Verbatim}[fontsize=\small,samepage=true,frame=single,framesep=3mm]
import datetime as dt
\end{Verbatim}

The package has these goodies:

\begin{itemize}
\itemsep1em

\index{timedelta@\texttt{timedelta} type (datetime)}

\item \textbf{The \texttt{timedelta} type}. Variables of this type represent
\textit{durations in time}. You can create them like this:

\begin{Verbatim}[fontsize=\small,samepage=true,frame=single,framesep=3mm]
one_semester = dt.timedelta(weeks=15)
one_year = dt.timedelta(weeks=52)
this_class = dt.timedelta(hours=1, minutes=15)
finger_snap = dt.timedelta(seconds=.1)
\end{Verbatim}

\vspace{-.15in}

You can also use \texttt{days}, \texttt{milliseconds}, and
\texttt{microseconds}. Anything else is ambiguous. (Ask yourself: why?)

\medskip

Don't think of the \texttt{timedelta} variable you create as being
intrinsically ``in'' any particular unit -- just think of it as an ``amount of
time,'' which could then be expressed in any units.

\texttt{timedelta}s can be positive or negative, by the way, and you can add or
subtract them to/from each other to get other \texttt{timedelta}s:

\begin{Verbatim}[fontsize=\small,samepage=true,frame=single,framesep=3mm]
print(one_semester - this_class)
print(this_class + finger_snap)
print(one_semester - one_year)
\end{Verbatim}
\vspace{-.2in}

\begin{Verbatim}[fontsize=\small,samepage=true,frame=leftline,framesep=5mm,framerule=1mm]
104 days, 22:45:00
1:15:00.100000
-259 days, 0:00:00
\end{Verbatim}

After printing the number of whole days (if any), the rest of the output gives
time in \textsl{hours}:\textsl{minutes}:\textsl{seconds}:\textsl{microseconds}. Note that the last of these three
outputs is negative, since \texttt{one\_semester} is obviously shorter than
\texttt{one\_year}.

\index{date@\texttt{date} type (datetime)}

\item \textbf{The \texttt{date} type}. Each \texttt{date} variable represents a
date in an idealized Gregorian calendar. You can create them via:
\begin{Verbatim}[fontsize=\small,samepage=true,frame=single,framesep=3mm]
election_day = dt.date(year=2020, month=11, day=3)
primary_day = dt.date(year=2020, month=3, day=4)
here_we_are = dt.date.today()
\end{Verbatim}

If you \textit{subtract} two \texttt{date}s, you get a \texttt{timedelta}. If
you \textit{add} a \texttt{date} to a \texttt{timedelta}, you get another
\texttt{date}.

\begin{Verbatim}[fontsize=\small,samepage=true,frame=single,framesep=3mm]
print(type(election_day - primary_day))
print(election_day - primary_day)
print(type(primary_day + one_semester))
print(primary_day + one_semester)
\end{Verbatim}

\begin{Verbatim}[fontsize=\small,samepage=true,frame=leftline,framesep=5mm,framerule=1mm]
datetime.timedelta
244 days, 0:00:00
datetime.date
2020-06-17
\end{Verbatim}

Note that if your \texttt{timedelta} has precision more granular than days, it
is lost in these operations:

\begin{Verbatim}[fontsize=\small,samepage=true,frame=single,framesep=3mm]
if dt.date.today() == dt.date.today() + finger_snap:
    print("It'll still be the same day a finger snap from now.")
if dt.date.today() == dt.date.today() + this_class:
    print("It'll still be the same day one lecture from now.")
if dt.date.today() == dt.date.today() + one_semester:
    print("It'll still be the same day one semester from now.")
\end{Verbatim}
\vspace{-.2in}

\begin{Verbatim}[fontsize=\small,samepage=true,frame=leftline,framesep=5mm,framerule=1mm]
It'll still be the same day a finger snap from now.
It'll still be the same day one lecture from now.
\end{Verbatim}
\normalsize

You can also compare dates with comparison operators like \texttt{<} and
\texttt{>=}:

\begin{Verbatim}[fontsize=\small,samepage=true,frame=single,framesep=3mm]
if primary_day < election_day:
    print("The primary occurs before the election, duh!")
else:
    print("Wuuuut??")
\end{Verbatim}
\vspace{-.2in}

\begin{Verbatim}[fontsize=\small,samepage=true,frame=leftline,framesep=5mm,framerule=1mm]
The primary occurs before the election, duh!
\end{Verbatim}

\index{time@\texttt{time} type (datetime)}

\item \textbf{The \texttt{time} type.} A ``time,'' by itself, represents an
abstract wall clock time, independent of any particular day. It doesn't mean
``5pm this afternoon'' but rather ``5pm's in general.''

\begin{Verbatim}[fontsize=\small,samepage=true,frame=single,framesep=3mm]
start_of_class = dt.time(hour=12, minute=30)
dinner_time = dt.time(hour=18, minute=0)
\end{Verbatim}

These can also be compared in expected ways, but \textit{not} added or
subtracted (which actually wouldn't make sense -- why?)

\index{datetime@\texttt{datetime} type (datetime)}

\item Finally, we have \textbf{the \texttt{datetime} type}, which as its name
suggests, combines features of both dates and times. In fact, a
\texttt{datetime} variable represents \textit{an absolute, specific point in
history}. Creation options include:

\begin{Verbatim}[fontsize=\small,samepage=true,frame=single,framesep=3mm]
this_exact_moment = dt.datetime.now()
terrorist_attack = dt.datetime(2001, 9, 11, 10, 33)
print("Right now it's: {}".format(this_exact_moment))
print("A horrible thing occurred at: {}".format(terrorist_attack))
\end{Verbatim}
\vspace{-.2in}

\begin{Verbatim}[fontsize=\small,samepage=true,frame=leftline,framesep=5mm,framerule=1mm]
Right now it's: 2020-03-10 09:33:57.015432
A horrible thing occurred at: 2001-09-11 10:33:00
\end{Verbatim}

\index{combine@\texttt{.combine()} method (\texttt{datetime})}

We can also use the \texttt{datetime}'s \texttt{.combine()} method to put a
\texttt{date} and a \texttt{time} together into a \texttt{datetime}:

\begin{Verbatim}[fontsize=\small,samepage=true,frame=single,framesep=3mm]
todays_class = dt.datetime.combine(here_we_are, start_of_class)
print("Today's lecture is at: {}".format(todays_class))
\end{Verbatim}
\vspace{-.2in}

\begin{Verbatim}[fontsize=\small,samepage=true,frame=leftline,framesep=5mm,framerule=1mm]
Today's lecture is at: 2020-03-10 12:30:00
\end{Verbatim}

\index{time@\texttt{.time()} method (of a \texttt{datetime} variable)}
\index{date@\texttt{.date()} method (of a \texttt{datetime} variable)}

\texttt{datetime}s can be subtracted from each other, and \texttt{timedelta}s
can be added or subtracted to/from them to produce other \texttt{datetime}s.
They can also be compared. Finally, \texttt{.time()} and \texttt{.date()}
methods will produce the constituent parts as \texttt{time} or \texttt{date}
variables, accordingly.

%\index{timestamp@\texttt{timestamp} type (datetime)}
%
%\small
%\item Oh, and for completeness, there's also a thing called a
%``\texttt{timestamp},'' which is really just a \texttt{float} value. A
%\texttt{timestamp} is -- get this -- \textit{the number of seconds that have
%elapsed since the stroke of midnight Jan.~1, 1970 UTC} (Coordinated Universal
%Time -- the Greenwich Mean Timezone).\footnote{That particular instant in time
%is sometimes called ``The Epoch.'' It's an arbitrary reference point, but was
%chosen in part to be approximately equal to Stephen's original birthday.} This
%will occasionally come up, since even though it's gibberish to humans, it is
%easy to work with mathematically.
%\normalsize

\end{itemize}


\section{Parsing and formatting}

Now most of the time we won't be creating \texttt{datetime} variables using the
techniques above. Instead, we'll be \textit{parsing} them from string data.

There are a seemingly unlimited number of ways that dates and times can be
expressed in strings. Just consider a few possibilities:

\small
\begin{compactitem}
\item Feb. 20, 2019 2pm
\item Feb 20, 2019 2:00pm
\item Feb 20, 2019 2:00 P.M.
\item 2/20/19 2:00 P.M.
\item 2/20/19 14:00
\item 02/20/19 14:00:00
\item 20-Feb-2019 2:00pm
\item 2019-Feb-20 14:00:00
\item ...
\end{compactitem}
\normalsize

and on it goes. It's a testimony to the power of the human mind that we can
look at these and at a glance infer what date/time was intended. Getting a
computer to automatically do so is somewhat of a challenge.

\subsection{Parsing with \texttt{.strptime()}}

The \texttt{date}, \texttt{time}, and \texttt{datetime} types all have a
(somewhat clunky) method called ``\texttt{.strptime()}'' that you can call,
passing (1) a string representing a date, time, or datetime, and (2) a
formatting string telling the method how to interpret it. This is best seen
through an example:

\begin{Verbatim}[fontsize=\small,samepage=true,frame=single]
my_string = "2/20/2019 14:00"
start_of_class = dt.datetime.strptime(my_string, "%m/%d/%Y %H:%M")

>>> print(start_of_class)
    2019-02-20 14:00:00
\end{Verbatim}

This example demonstrates a full ``round-trip'' operation: we started with a
\texttt{string}, parsed it into a \texttt{datetime}, and then printed it back
to the screen as a \texttt{string}. Note that the format of the printout
didn't match the format of ``\texttt{my\_string}''; this is okay. What matters
internally to our program is that our date \& time variables \textit{represent}
the right things, not that they are printed in certain ways.

Now what are all those funky percent signs and junk in the second argument to
\texttt{.strptime()}? They all have a special meaning. Here are some common
ones (others are in the docs):

\begin{center}
\begin{tabular}{|c|l|}
\hline
\%a & Weekday (abbreviated): Sun, Mon, \textit{etc.} \\
\%A & Weekday (full): Sunday, Monday, \textit{etc.} \\
\%w & Weekday (number): 0, 1, ..., 6 \\
\%d & Day of month (number): 0, 1, ..., 31 \\
\%b & Month (abbreviated): Jan, Feb, \textit{etc.} \\
\%B & Month (full): January, February, \textit{etc.} \\
\%m & Month (number): 01, 02, ..., 12 \\
\%y & \textbf{Two}-digit year: 00, 01, ..., 99 \\
\%Y & \textbf{Four}-digit year: ..., 1999, 2000, 2001, ...\\
\hline
\%H & Hour on \textbf{24-hour} clock: 01, 02, ..., 23 \\
\%I & Hour on \textbf{12-hour} clock: 01, 02, ..., 12 \\
\%p & AM/PM (only used in conjunction with \textbf{\%I}) \\
\%M & Minute: 01, 02, ..., 59 \\
\%S & Second: 01, 02, ..., 59 \\
\hline
\end{tabular}
\end{center}

\pagebreak
You have to be really careful here and just go slow. Examine the date and/or
time input that you have in your data set, and work out which bits of info,
combined with which literal characters (like ``\texttt{/}'' and ``\texttt{:}''
in the example above), appear in which order.

\subsection{Formatting with \texttt{.strftime()}}

Complementarily to \texttt{.strptime()} is ``\texttt{.strftime()}''. (Note that
the ``p'' stands for ``\textbf{p}arse,'' while the ``f'' stands for
``\textbf{f}ormat.'') This is the reverse operation: given a \texttt{date},
\texttt{time}, or \texttt{datetime} variable, spit it out as a string of
characters. This isn't used nearly as often, but it occasionally is, and it's
useful to know as a translation mechanism.

\begin{Verbatim}[fontsize=\small,samepage=true,frame=single]
>>> print(start_of_class.strftime("%A, %B the %dth at %I%p."))
    Tuesday, February the 20th at 02PM.
\end{Verbatim}

One common error: note that you call \texttt{.strptime()} directly on the
\textit{type} (\texttt{date}, \texttt{time}, or \texttt{datetime}), passing
two arguments:

\begin{Verbatim}[fontsize=\small,samepage=true,frame=single]
an_object = dt.datetime.strptime(my_birthday_string, "%m/%d/%Y %H:%M")
\end{Verbatim}

but you call \texttt{.strftime()} on the object in question
(a \texttt{date}, \texttt{time}, or \texttt{datetime}), passing
one argument:

\begin{Verbatim}[fontsize=\small,samepage=true,frame=single]
a_string = my_birthday_object.strftime("%m/%d/%Y %H:%M")
\end{Verbatim}


